\chapter*{Abstract}
\sloppy

\vspace{1cm}
\par Government Identity documents have become groundwork for citizen verification, financial inclusion and public service. However unauthorized disclosure, fraudulent access and frequent misuse of individual's personal information expose gaps in routine verification and wrongful disbursement of welfare benefits which asks the immediate need for a more secure privacy holding approach. Existing infrastructure like DigiLocker takes care of the secure transmission but a factor of privacy exposure is still compromised. The proposed model introduces TrustChain based on blockchain framework for decentralized identify verification and secure access. The inducement shifts focus from submission of identity information to authentication by the means of DID (Decentralized Identifiers) ensuring that some personal information will be hidden to the document requestor. By integrating self sovereign identity principles, distributed storage and cryptographic operations the model enables users to authenticate without revealing personal parameters minimizing the risks of identity compromise. Pilot findings are particular not to mention that identity exposure is mitigated by such a representation and offers scalability towards integration into the current infrastructures such as Aadhaar connected services and Digital locker. A safe identity space where individuals retain ownership to personal information as organizations and governments achieve acceptable validation. 

\paragraph{}





\vspace*{0.7cm}
%\vspace*{1cm}
\noindent
{\large{\textbf{\emph{Keywords:}}}}~
Blockchain Framework, DID (Decentralized Identifiers), Self Sovereign Identity, Privacy Preserving Verification, Zero Knowledge Proofs, Secure Access





